\section{引言}\label{sec:introduction}

此处填入按原文顺序翻译的引言。第一次出现的术语建议写成中文名称（English Term，ET）；正文应保留作者对研究结果的限定条件与引用位置。此处演示可点击的引用\cite{sample-entry}。

如下为演示公式，而非原论文的实验公式：
\begin{equation}
    x_{k+1}=x_k+u_k.
    \label{eq:illustration}
\end{equation}

式~\eqref{eq:illustration} 的编号由 \LaTeX{} 统一管理。正式译排时保留原式编号；编号不连续时可使用 \verb|\tag{...}|，同时配置语义化的 \verb|\label|。

% Example for a real vector figure (replace with extracted original figure):
% \begin{figure}[htbp]
%   \centering
%   \includegraphics[width=.94\linewidth]{figures/fig01.pdf}
%   \caption{按原意翻译的图注。}
%   \label{fig:architecture}
% \end{figure}
